% Capítulo 3: Marco Teórico / Estado del Arte

\chapter{Marco Teórico / Estado del Arte}

\label{Chapter3}

\lhead{Capítulo 3. \emph{Marco Teórico}}

\textit{[En un párrafo explica brevemente los contenidos de este capítulo antes de la distribución en distintos apartados. En un marco teórico se debe describir la información científica relevante para tu proyecto. Procura reflejar información actualizada y citarla adecuadamente (normas APA, 6ª edición). Incluye al menos tantos apartados como áreas en las que esté implicado tu proyecto.]}

\textit{[Fuentes habituales donde obtener trabajos de referencia: Google Académico (\url{https://scholar.google.es/}), ERIC (\url{https://eric.ed.gov/}), Psycnet (\url{http://psycnet.apa.org/}), Dialnet (\url{https://dialnet.unirioja.es/}), Medline/PubMed (\url{https://www.nlm.nih.gov/bsd/pmresources.html}).]}

\textit{[Se recomienda un tamaño de 10 a 20 páginas para este capítulo.]}

%----------------------------------------------------------------------------------------

\section{}

%----------------------------------------------------------------------------------------

\section{}

%----------------------------------------------------------------------------------------

\section{}
